
Орчин үеийн программ хангамжийн хөгжүүлэлтэд framework-ийн үүрэг улам бүр нэмэгдэж, ялангуяа backend систем боловсруулах үйл явцад чухал байр суурь эзэлж байна.Backend систем хөгжүүлэхэд давтагддаг ажлуудыг автоматжуулах, хөгжүүлэгчийн бүтээмжийг нэмэгдүүлэх, кодын хэмжээг багасгах, мөн системийн уян хатан болон өргөжих чадварыг сайжруулах шаардлага улам бүр өсөж байна.

Java хэл дээр хөгжүүлсэн \textbf{Spring Framework} нь annotation-д суурилсан тохиргоо (\textit{annotation-based configuration}) болон \textbf{Dependency Injection (DI)} зарчмыг нэвтрүүлснээр дэлхий даяар хамгийн өргөн хэрэглэгддэг backend framework-үүдийн нэг болсон. Тус framework нь \textbf{Inversion of Control (IoC)} болон \textbf{Dependency Injection (DI)} зэрэг архитектурын оновчтой шийдлүүдийг ашигласнаар хөгжүүлэгчийн бичих шаардлагатай кодын хэмжээг эрс бууруулж, системийн өргөжих чадвар (\textit{scalability}), засварлах болон удирдах боломжийг мэдэгдэхүйц нэмэгдүүлдэг.

Гэсэн хэдий ч ихэнх хөгжүүлэгчид эдгээр framework-үүдийг хэрэглээний түвшинд ашигладаг боловч тэдгээрийн дотоод ажиллагааны үндэс болох \textbf{Java Reflection API}, annotation боловсруулах механизм, объектын удирдлагын мөчлөг зэрэг суурь ойлголтуудыг төдийлөн гүнзгий судалдаггүй. Үүний улмаас framework-ийн дотоод архитектур, ажиллах зарчим, өргөтгөх боломжийг бүрэн ойлгохгүй байх тохиолдол түгээмэл ажиглагддаг.

Иймээс backend системийн суурь архитектур, объектын амьдралын мөчлөгийн удирдлага, annotation болон reflection-д суурилсан илрүүлэлт, удирдлагын аргачлалыг онолын болон практик түвшинд судлах шаардлага бий болж байна. Энэ нь framework-ийн цөм зарчим болох \textbf{IoC}, \textbf{DI}, \textbf{annotation} боловсруулалт, \textbf{reflection} ашиглалт зэрэг ойлголтуудыг илүү гүнзгий эзэмшихээс гадна backend хөгжүүлэлтийн суурь архитектурын мэдлэгийг бататгах, цаашдын судалгаа болон практик хэрэглээнд ач холбогдолтой юм.

\subfile{goal}